%File: anonymous-submission-latex-2023.tex
\documentclass[letterpaper]{article} % DO NOT CHANGE THIS
\usepackage{aaai23}  % DO NOT CHANGE THIS
\usepackage{times}  % DO NOT CHANGE THIS
\usepackage{helvet}  % DO NOT CHANGE THIS
\usepackage{courier}  % DO NOT CHANGE THIS
\usepackage[hyphens]{url}  % DO NOT CHANGE THIS
\usepackage{graphicx} % DO NOT CHANGE THIS
\urlstyle{rm} % DO NOT CHANGE THIS
\def\UrlFont{\rm}  % DO NOT CHANGE THIS
\usepackage{natbib}  % DO NOT CHANGE THIS AND DO NOT ADD ANY OPTIONS TO IT
\usepackage{caption} % DO NOT CHANGE THIS AND DO NOT ADD ANY OPTIONS TO IT

\frenchspacing  % DO NOT CHANGE THIS
\setlength{\pdfpagewidth}{8.5in} % DO NOT CHANGE THIS
\setlength{\pdfpageheight}{11in} % DO NOT CHANGE THIS
%
% These are recommended to typeset algorithms but not required. See the subsubsection on algorithms. Remove them if you don't have algorithms in your paper.
\usepackage{algorithm}
\usepackage{algorithmic}


%%%%% ADDED PACKAGES

\usepackage{subcaption}
\usepackage{paralist}
%
% These are are recommended to typeset listings but not required. See the subsubsection on listing. Remove this block if you don't have listings in your paper.
\usepackage{newfloat}
\usepackage{listings}
\DeclareCaptionStyle{ruled}{labelfont=normalfont,labelsep=colon,strut=off} % DO NOT CHANGE THIS
\lstset{%
	basicstyle={\footnotesize\ttfamily},% footnotesize acceptable for monospace
	numbers=left,numberstyle=\footnotesize,xleftmargin=2em,% show line numbers, remove this entire line if you don't want the numbers.
	aboveskip=0pt,belowskip=0pt,%
	showstringspaces=false,tabsize=2,breaklines=true}
\floatstyle{ruled}
\newfloat{listing}{tb}{lst}{}
\floatname{listing}{Listing}
%
% Keep the \pdfinfo as shown here. There's no need
% for you to add the /Title and /Author tags.
\pdfinfo{
/TemplateVersion (2023.1)
}

% DISALLOWED PACKAGES
% \usepackage{authblk} -- This package is specifically forbidden
% \usepackage{balance} -- This package is specifically forbidden
% \usepackage{color (if used in text)
% \usepackage{CJK} -- This package is specifically forbidden
% \usepackage{float} -- This package is specifically forbidden
% \usepackage{flushend} -- This package is specifically forbidden
% \usepackage{fontenc} -- This package is specifically forbidden
% \usepackage{fullpage} -- This package is specifically forbidden
% \usepackage{geometry} -- This package is specifically forbidden
% \usepackage{grffile} -- This package is specifically forbidden
% \usepackage{hyperref} -- This package is specifically forbidden
% \usepackage{navigator} -- This package is specifically forbidden
% (or any other package that embeds links such as navigator or hyperref)
% \indentfirst} -- This package is specifically forbidden
% \layout} -- This package is specifically forbidden
% \multicol} -- This package is specifically forbidden
% \nameref} -- This package is specifically forbidden
% \usepackage{savetrees} -- This package is specifically forbidden
% \usepackage{setspace} -- This package is specifically forbidden
% \usepackage{stfloats} -- This package is specifically forbidden
% \usepackage{tabu} -- This package is specifically forbidden
% \usepackage{titlesec} -- This package is specifically forbidden
% \usepackage{tocbibind} -- This package is specifically forbidden
% \usepackage{ulem} -- This package is specifically forbidden
% \usepackage{wrapfig} -- This package is specifically forbidden
% DISALLOWED COMMANDS
% \nocopyright -- Your paper will not be published if you use this command
% \addtolength -- This command may not be used
% \balance -- This command may not be used
% \baselinestretch -- Your paper will not be published if you use this command
% \clearpage -- No page breaks of any kind may be used for the final version of your paper
% \columnsep -- This command may not be used
% \newpage -- No page breaks of any kind may be used for the final version of your paper
% \pagebreak -- No page breaks of any kind may be used for the final version of your paper
% \pagestyle -- This command may not be used
% \tiny -- This is not an acceptable font size.
% \vspace{- -- No negative value may be used in proximity of a caption, figure, table, section, subsection, subsubsection, or reference
% \vskip{- -- No negative value may be used to alter spacing above or below a caption, figure, table, section, subsection, subsubsection, or reference

\setcounter{secnumdepth}{0} %May be changed to 1 or 2 if section numbers are desired.

% The file aaai23.sty is the style file for AAAI Press
% proceedings, working notes, and technical reports.
%

% Title

% Your title must be in mixed case, not sentence case.
% That means all verbs (including short verbs like be, is, using, and go),
% nouns, adverbs, and adjectives should be capitalized, including both words in hyphenated terms, while
% articles, conjunctions, and prepositions are lower case unless they
% directly follow a colon or long dash

\title{Towards integrated learning and constraint reasoning}
\author{
    %Authors
    % All authors must be in the same font size and format.
    Emilio Gamba\textsuperscript{\rm1, \rm 2},
    Tias Guns\textsuperscript{\rm1, \rm 2}
%    Maxime Mulamba\textsuperscript{\rm1, \rm 2},
%    Ignace Bleukx\textsuperscript{\rm 2},
%    Senne Berden\textsuperscript{\rm 2},
%    Milan Pesa\textsuperscript{\rm 2},
}
\affiliations{
    %Afiliations
    \textsuperscript{\rm 1}Vrije Universiteit Brussel, 
%    Pleinlaan 2, 1050 Brussels, Belgium \\
%    \{emilio.gamba, maxime.mulamba\}@vub.be,\\
\{firstname.lastname@vub.be\} \\
    \textsuperscript{\rm 2}KU Leuven, 
%    Oude Markt 13, 3000 Leuven, Belgium \\
    % email address must be in roman text type, not monospace or sans serif
%    \textsuperscript{\rm 1} firstname.lastname@vub.be, 
%\{tias.guns, ignace.bleuckx, senne.berden, milan.pesa\}@kuleuven.be %\\
\{firstname.lastname\}@kuleuven.be \\
    %\textsuperscript{\rm 3}
%
% See more examples next
}


\begin{document}

\maketitle
 
In recent years, the combination of deep learning for image perception and symbolic reasoning to visual sudoku has become increasingly popular.
Some systems jointly learn the constraints to classify cell values~\cite{Pushing, SATNet, NeurASP}. 
However this requires solver-specific mechanisms whereas our approach works with any constraint  solver. 
Another limitation, also present in \cite{CLRD, mulamba2020hybrid} is that instances are built by sampling images from MNIST \cite{mnist}, on which recent machine learning approaches achieve near-perfect accuracy \cite{ensemble-mnist}. Images from a phone camera are in RGB space, can contain artefacts such as grid borders and are overall noisier depending on the angle of the camera or lighting conditions. They are a bigger challenge. 

\section{Description of the Tutorial}


Constraint solving is a well-established approach in AI for reasoning. It involves formulating and solving complex constraint satisfaction and optimisation problems. 
In this tutorial, we take a look at different approaches to interpret and explain a pen-and-paper Sudoku scanned with a smartphone.

%% B6: Constraint Programming and Machine Learning
%Bringing together CP (Constraint Programming) and ML (Machine Learning) is an important aspect of the larger goal of integrating Reasoning and Learning. Participants are not expected to have prior experience in both fields, but to have familiarity with each at least at the level of an introductory AI course.
%The Bridge is designed to provide attendees with a better, broader sense of where we are and where we should be going, and an opportunity to brainstorm, discuss, debate, and find collaborators. The focus will be on bringing together the traditional AI fields of constraint-based reasoning and machine learning, but participants from related fields of reasoning, optimization, and learning, e.g., SAT, operations research, data mining, will be welcome. It is hoped that this one-day Bridge event will help establish an ongoing community, leading to workshops, special issues, etc.

%Where Are We Now?
%o Mini-tutorials or mini-surveys. They could be broad or narrow in scope, restricted to a particular subtopic, covering recent work from a particular conference or workshop series, or even introducing or surveying a body of the presenters’ work.

%Where Do We Go from Here?
%o Position papers, grand challenges, proposals, roadmaps, questions, needs, objectives, opportunities, obstacles, testbeds, applications, metrics. Proposals for discussion topics, panels, or debates. Introductions to ongoing research projects.

%Pitches
%o Presentations by those seeking collaborators or mentors for current or proposed work.

%Tools
%o Mini-surveys of available CP or ML tools, or mini-tutorials on individual tools. Examples of the use of CP tools for ML or vice versa. Also welcome are proposals for “hands on” sessions where participants can be guided through a real-time introduction to a tool they access on their laptops.

%Community Building
%o Proposals for special issues, conference tracks, workshops, competitions, tools, libraries, benchmark problems, websites, tutorials, surveys, talks, funding proposals, collaborations, blogs, groups, videos, community outreach, social media, instructional material, syllabi, ontologies, books, apps, testbeds, bibliographies, etc.

\section{Proposed format of the Tutorial}
The tutorial should last around 1:30.

\subsection{Detailed, point-form outline of the tutorial}
\begin{enumerate}
	\item Quick introduction to Constraint Solving 
	\begin{itemize}
		\item Examples demonstrating the language constructs
	\end{itemize}
	\item Combining prediction and reasoning
	\begin{itemize}
		\item State-of-the-art techniques for combining CP with ML
		\item Integrating the prediction (ML) into reasoning (CP)
	\end{itemize}
	\item Use-cases: Where to go from there ? 
	\begin{itemize}
		\item Explanations 
		\item ...
	\end{itemize}
	\item Conclusion, Outlooks and Questions
\end{enumerate}

\subsection{Discussion starters}

%Given the condensed time frame, submissions need only be one- or two-page abstracts, outlining the proposed presentation, the minimum and maximum time it could occupy, and citing relevant experience/publications of the author(s).
%Proposals can be for presentations, or portions of presentations, that have been made elsewhere (e.g. tutorials or surveys), suitably updated and scaled for this Bridge event. Presentations (aside from Tools) need to combine CP and ML in some fashion. All presentations should keep in mind that audience members are not expected to have expertise in both areas. Presentations should “reach across the aisle”, and welcome questions and discussion.

%\section{Hybrid prediction and reasoning}
%
%\section{Related Works}
%
%\section{Conclusion}

\clearpage
\section{Acknowledgements}
This research received partial funding from the Flemish Government (AI Research Program); and funding from the European Research Council (ERC) under the European Union’s Horizon 2020 research and innovation program (Grant No. 101002802, CHAT-Opt).


\bibliography{aaai23.bib}  
% \nobibliography{aaai23}


\end{document}
